\documentclass[10pt,a4paper]{amsart}
\usepackage[margin=19mm]{geometry}
\usepackage{amsmath,amssymb,mathtools}
\usepackage[scr=rsfs,cal=euler]{mathalfa}
\usepackage{xfrac}
\usepackage[unicode,hidelinks,bookmarks=false]{hyperref}
\newcommand{\ie}{i.e.\ }
\newcommand{\cf}{cf.\ }
\newcommand{\resp}{respectively\ }
\newcommand{\Subsection}{\subsection}
\providecommand{\nobreakdash}{\nobreak\hbox{-}}
\setlength{\emergencystretch}{2em}
\multlinegap0pt
\usepackage{amssymb}
\usepackage[shortlabels]{enumitem}
\usepackage{xcolor}
\usepackage{bm}
\usepackage{mathtools}
\newtheorem{lemmaen}{Lemma}
\newtheorem{proposition}{Proposition}
\usepackage{cleveref}
\crefname{lemmaen}{Lemma}{Lemmas}
\crefname{proposition}{Proposition}{Propositions}
\newcommand{\R}{\mathbb{R}}
\newcommand{\de}{\partial}
\DeclareMathOperator{\dive}{div}
\newcommand{\hau}{\mathcal{H}}
\newcommand{\str}{\mathcal{U}}
\title{Michael--Simon inequality \\ for anisotropic energies close to the area \\ via multilinear Kakeya-type bounds}
\author{Guido De Philippis, Alessandro Pigati}
\date{Source: arXiv:2601.10647v3 (CC BY 4.0)}
\begin{document}
\maketitle

\noindent\emph{Source and licence.} Michael--Simon inequality \\ for anisotropic energies close to the area \\ via multilinear Kakeya-type bounds. Version arXiv:2601.10647v3 (CC BY 4.0), distributed by arXiv under the Creative Commons Attribution 4.0 International licence, http://creativecommons.org/licenses/by/4.0/, as recorded in the arXiv licence metadata for this exact version and shown on the abstract page https://arxiv.org/abs/2601.10647v3. Full licence notice: \texttt{licenses/arxiv-2601.10647-CC-BY-4.0.txt}.\par\smallskip

\noindent\textit{Editorial scope.} Counted unit: the article's proposition complement.bis (source0.tex lines 1208--1220). The article's complete original proof of it is delivered with it (source0.tex lines 1222--1268, one \texttt{proof} environment of 47 lines). No mathematics is written, repaired or re-derived: the statement and the proof are the article's own lines, in the article's own notation, and no retained line is substituted or re-flowed.  Retained uncounted as the source-local context the proof needs: lines 1131--1148.  Nothing was omitted from any retained span, so the package carries no omission record.  No retained line contains a citation, so the wrapper carries no bibliography and no bibliography entry.  No retained line contains a figure environment, an image-inclusion command or drawing code, so no image is delivered and the package declares no asset dependency.  Editorial formatting only, in addition: an AMS \texttt{amsart} preamble that restates the article's own declarations and macros used by the retained lines, namely the environments \texttt{lemmaen}, \texttt{proposition} on separate counters, together with the standard packages those lines need.  The retained environments print in this document as Lemmaen 1, Proposition 1 in order of appearance, whereas the article prints the same items with the numbering of its own full document.  The complete native source of the article as distributed in the arXiv e-print for this exact version is bundled unchanged as \texttt{sources/192787/source0.tex}.
\begin{lemmaen}\label{lemma.curve}
%There exist two universal constants $\Lambda,C'\ge1$ such that
The following hold
for a Lipschitz curve $\gamma:[0,L]\to\str_j^x=[j,j+1]\times\R$ with unit speed:
\begin{itemize}
\item[(i)] if %$\gamma$ is a loop, or alternatively $\{\gamma^x(0),\gamma^x(L)\}\subseteq\{j,j+1\}$ and
$\gamma^x(0)\ge\gamma^x(L)$
then
$$\int_0^L(\dot\gamma^x-|\dot\gamma^y|)^+\,dt
\le \int_0^L(\dot\gamma^x-|\dot\gamma^y|)^-\,dt;$$
\item[(ii)] %if either endpoint of $\gamma$ belongs to the open stripe $(j,j+1)\times\R$,
in general, without the previous assumption, we have
$$\int_0^L(\dot\gamma^x-|\dot\gamma^y|)^+\,dt
\le 1+\int_0^L(\dot\gamma^x-|\dot\gamma^y|)^-\,dt.$$
%If $\dot\gamma^x\ge0$, then $C'$ can be taken arbitrarily small (by increasing $\Lambda$).
\end{itemize}
Analogous statements hold for a stripe $\str_k^y$, interchanging the roles of $x,y$.
\end{lemmaen}

\begin{proposition}\label{complement.bis}
	We have
	\begin{align*}
		&\int_{\mathcal{R}^x}(S_d^x-|S_d^y|)^+
		\le\int_{\mathcal{R}^x}\lvert \dive S\rvert+\int_{\str_j^x}(S_d^x-|S_d^y|)^-.
	\end{align*}
    %If $S^x\ge0$ then we can take $C'=1$.
%	As a consequence,
%	\begin{align*}
%		&\int_{\str_j^x}|S_d|
%		\le\Lambda\int_{\str_j^x}\lvert \dive S\rvert+\frac{1}{\Lambda-1}\int_{\str_j^x}(S_d^x-|S_d^y|)^-.
%	\end{align*}
\end{proposition}

\begin{proof}
	Writing $Z=\nabla^\perp f$, consider a level set $\gamma:[0,L]\to\mathcal{R}^x$ of $f$, parametrized by arclength.
	We have
	\begin{align*}
		&\int_0^L\alpha_d(\dot\gamma^x-|\dot\gamma^y|)^+\,dt
		=\int_0^{\max\alpha_d}\int_{\{\alpha_d>s\}}(\dot\gamma^x-|\dot\gamma^y|)^+\,dt\,ds,
	\end{align*}
	where we write $\alpha_d$ in place of $\alpha_d\circ\gamma$ for simplicity.
	Note that, calling $N_s$ the cardinality of $\{\alpha_d=s\}$, the set $\{\alpha_d>s\}$
	has at most $N_s$ connected components, since $\alpha_d$
	takes all values in $(0,\max\alpha_d)$. Hence, applying part (ii) of
	\cref{lemma.curve} to $\gamma$ restricted to each connected component, we get
	$$\int_{\{\alpha_d>s\}}(\dot\gamma^x-|\dot\gamma^y|)^+\,dt
	\le N_s+\int_{\{\alpha_d>s\}}(\dot\gamma^x-|\dot\gamma^y|)^-\,dt.$$
	Integrating in $s$ and using the area formula, we deduce that
	$$\int_0^L\alpha_d(\dot\gamma^x-|\dot\gamma^y|)^+
	\le \int_0^L|\dot\alpha_d|+\int_0^L\alpha_d(\dot\gamma^x-|\dot\gamma^y|)^-.$$
	%We distinguish two cases: if the length $\hau^1(\gamma\cap\{\alpha_d>s\})\le\Lambda$, we just bound the inner integral by $\Lambda$.
	%
	%Otherwise, note that, since $\dot\gamma^x+|\dot\gamma^y|\ge1$ and $\int_{\{\alpha_d>s\}}\dot\gamma^x\le\int_0^L\dot\gamma^x=1$ (as $\dot\gamma^x>0$), we must have
	%\begin{align*}
	%	&\int_{\{\alpha_d>s\}}(\dot\gamma^x-|\dot\gamma^y|)^-
	%	\ge\int_{\{\alpha_d>s\}}(|\dot\gamma^y|-\dot\gamma^x)
	%	\ge\int_{\{\alpha_d>s\}}(1-2\dot\gamma^x)
	%	\ge \Lambda-2,
	%\end{align*}
	%which gives
	%\begin{align*}
	%	&\int_{\{\alpha_d>s\}}(\dot\gamma^x-|\dot\gamma^y|)^+
	%	\le\int_{\{\alpha_d>s\}}\dot\gamma^x
	%	\le 1
	%	\le\frac{1}{\Lambda-2}\int_{\{\alpha_d>s\}}(\dot\gamma^x-|\dot\gamma^y|)^-.
	%\end{align*}
	%
	%Also, the maximum of $\alpha_d$ along $\gamma$ is bounded by $\int_0^L|\dot\alpha_d|$.
%If $S^x\ge0$ then $C'$ can be taken as small as desired, as $\dot\gamma^x\ge0$
%in this case; in particular, we can take $C'=1$ (with $\Lambda:=3$).
    
	Since $\dot\gamma=\frac{Z}{|Z|}$ and $\dive Z=0$, we have $|\dot\alpha_d|=\frac{\lvert \dive S_d\rvert}{|Z|}=\frac{\lvert \dive S\rvert}{|df|}$, where we omit composition with $\gamma$. To sum up, recalling that $\dot\gamma=\frac{(-\de_y f,\de_x f)}{|df|}$, we get
	\begin{align*}
		\int_0^L\alpha_d\frac{(-\de_y f-|\de_x f|)^+}{|df|}\circ\gamma
		%&\le\Lambda\max\alpha_d+\frac{1}{\Lambda-2}\int_0^\infty\int_{\{\alpha_d>s\}}\frac{(-\de_y f-|\de_x f|)^-}{|df|}\circ\gamma(t)\,dt\,ds \\
		&\le\int_0^L\frac{\lvert \dive S\rvert}{|df|}\circ\gamma
		+\int_0^L\alpha_d\frac{(-\de_y f-|\de_x f|)^-}{|df|}\circ\gamma.
	\end{align*}
	The claim follows using the coarea formula for $f$.
\end{proof}

\end{document}
